\documentclass[11pt]{article}
\usepackage{fontspec}
\usepackage[french,english]{babel}
\usepackage{amsmath}
\usepackage{amssymb}
\usepackage{amsthm}
\usepackage{mathrsfs}
\usepackage{tikz-cd}
\usepackage{enumitem}
\usetikzlibrary{babel}

\newcommand{\printedpage}[1]{\marginpar{\footnotesize p.~#1}}
\newcommand{\msp}[1]{\marginpar{\footnotesize [#1]}}
\newcommand{\cM}{\mathcal{M}}
\newcommand{\HO}{\mathrm{HO}}
\newcommand{\Hom}{\mathrm{Hom}}
\newcommand{\Ob}{\mathrm{Ob}}
\newcommand{\rf}{\mathrm{f}}
\newcommand{\rc}{\mathrm{c}}

\begin{document}

\printedpage{34}
% head: A. Grothendieck : Dérivateurs, Ch. XIII, Édition des Universités de Montpellier II et Paris VII

\begin{itemize}[leftmargin=3em]
\item[c)] \textit{In particular, if $I'$, $I$ are finite ordered sets, then every $u : I' \longrightarrow I$ is admissible. The same holds, more generally, if $I$, $I'$ are ordered sets and $I'^{\circ}$ is a noetherian sober space.}
\end{itemize}

\begin{itemize}[leftmargin=2em]
\item[2)] \textit{Suppose $u : I' \longrightarrow I$ admissible. Consider}
\[
u_*^W \ : \ \underbrace{\cM(I')_{\rf}W(I')_{\rf}^{-1}}_{\simeq\, \HO_W(I')} \ \longrightarrow \ \underbrace{\cM(I)_{\rf}W(I)_{\rf}^{-1}}_{\HO_W(I)},
\]
\textit{induced by $u_*$. This functor is right adjoint to $u_W^* : \HO_W(I) \longrightarrow \HO_W(I')$.}

\msp{page 61}
\textit{Moreover, the canonical homomorphism (for $s \in \Ob I$, $X' \in \Ob \HO_W(I')$)}
\[
u_*^W(X')(s) \ \simeq \ p_{s\,*}(X'|I'/s) \ \overset{\text{def}}{=} \ \underset{I'/s}{\underleftarrow{\mathrm{holim}}}{}_W(X'|I'/s)
\]
\textit{(where $p_s : I'/s \longrightarrow e$) is an isomorphism, if we suppose that the $I'/s$ are admissible (we shall then say that $u$ is} strictly \textit{admissible).}
\end{itemize}

\textbf{Corollary.} \textit{On the domain formed by finite ordered sets, the prederivator} $\HO_W$ \textit{satisfies (in addition to axioms} \textbf{Der}~1, \textbf{Der}~2\textit{) axioms} \textbf{Der}~4 \textit{and} \textbf{Der}~5 \textit{(existence of the direct images $u_*$, right adjoints of the functors $u^*$, and standard calculation of the fibers of a direct image).}

\textsc{Proof} of the theorem-scholie.

$1^\circ$) \textit{a}) is immediate. For \textit{b}), we proceed by induction on the cardinality $n$ of the finite ordered set $I$, the cases $n \in \{0,1\}$ being trivial. Suppose then $n \geq 2$, and the theorem proved for the $n' < n$. I distinguish two cases: if $I$ has a greatest element $a$, then, denoting by $\alpha : e \longrightarrow I$ the functor with value $a$, by $p : I \longrightarrow e$ the projection, we know that we have
\[
p_* \ \simeq \ \alpha^*, \qquad \textit{i.e.} \ \varprojlim_{I} X_i \simeq X_a,
\]
whence the fact that on the whole of $\cM(I)$
\msp{page 62}
(not only on $\cM(I)_{\rf}$), $p_*$ is compatible with the localizers $W(I)$, $W$.

If $I$ has no greatest element, then the set $M$ of its maximal elements has cardinality $\geq 2$, we can therefore write it $M = M_1 \amalg M_2$, and we will have % [sic] des ses -> de ses, typographical slip
\[
I \ = \ I_1 \cup I_2, \qquad \text{where } I_1 = I/M_1,\ I_2 = I/M_2
\]
are two open subsets of $I$, different from $I$, hence of cardinalities $n_1$, $n_2 \ < \ n$. Consider also $I_0 = I_1 \cap I_2$, it is an open subset of cardinality $< n$. The diagram
\[
\begin{tikzcd}[row sep=small]
 & I_1 \arrow[dr, dotted] & \\
I_0 \arrow[ur] \arrow[dr] & & I = I_1 \cup I_2 \\
 & I_2 \arrow[ur, dotted] &
\end{tikzcd}
\]
% foot: Version 25 janvier 2004
\printedpage{35}
% head: A. Grothendieck : Dérivateurs, Ch. XIII, Édition des Universités de Montpellier II et Paris VII
corresponds to an arrow
\[
\begin{tikzcd}[row sep=small, column sep=small]
 & & & 1 & \\
u \ : \ I \arrow[r] & \Psi = 0 \arrow[urr, start anchor=east, end anchor=south west] \arrow[drr, start anchor=east, end anchor=north west] & & & , \\
 & & & 2 &
\end{tikzcd}
\]
with
\[
I_1 \ = \ I/1, \qquad I_2 \ = \ I/2, \qquad I_0 \ = \ I/0.
\]
The induction hypothesis and \textit{a}) imply that $u$ is admissible. As a composite of admissible functors is admissible (tautologically, via $(vu)_* \simeq v_*u_*$), we are finally reduced to the
\msp{page 63}

\textbf{Lemma.} \textit{The category} \begin{tikzcd}[row sep=small, column sep=small, baseline=(\tikzcdmatrixname-2-1.base)]
 & 1 \\
\Psi = 0 \arrow[ur] \arrow[dr] & \\
 & 2
\end{tikzcd} \quad \textit{is admissible.}

This will follow from the

\textbf{Corollary} (\footnote[36]{\textbf{NB} There is an obvious variant, characterizing the $X$ such that $X \longrightarrow e$ is in $TF(\Psi)$ as those for which $X_0$, $X_1$, $X_2$ are fibrant and \textit{trivial,} and $\alpha_1$, $\alpha_2 \in TF$.}). \textit{Let} \begin{tikzcd}[row sep=small, column sep=small, baseline=(\tikzcdmatrixname-2-1.base)]
 & X_1 \arrow[dl, "\alpha_1"'] \\
X_\bullet = X_0 & \\
 & X_2 \arrow[ul, "\alpha_2"]
\end{tikzcd} \quad \textit{be in $\cM(\Psi)$. For $X_\bullet$ to be fibrant, it is necessary and sufficient that we have}

\begin{itemize}[leftmargin=3em]
\item[a)] \textit{$X_0$, $X_1$, $X_2$ are fibrant, and}
\item[b)] \textit{$\alpha_1$, $\alpha_2 \in F$.}
\end{itemize}

The necessity of \textit{a}) is a special case of corollary 2, page 48, that of \textit{b}) is a special case of corollary 1 page 50. For sufficiency (which we shall not need for the theorem), we know how to show that in the commutative diagram in solid lines
\[
% Audit resolved: PDF leaf 7 confirms interrupted vertical sides A_0 -> X_0 and A_2 -> X_2, obscured by the displaced B labels; the printed geometry is retained. See corrections.csv.
\begin{tikzcd}
A_1 \arrow[r] \arrow[dr, "TC"] \arrow[dd] & A_0 \arrow[dr, "TC"] \arrow[d, dash] & A_2 \arrow[l] \arrow[dr, "TC"] \arrow[d, dash] & \\
 & B_1 \arrow[r] \arrow[dl, dotted, "h_1"'] \arrow[d] & B_0 \arrow[dl, dotted, "h_0"'] \arrow[d] & B_2 \arrow[l] \arrow[dl, dotted, "h_2"'] \\
X_1 \arrow[r, "\alpha_1"'] & X_0 & X_2 \ , \arrow[l, "\alpha_2"] &
\end{tikzcd}
\]
with three arrows marked $TC$, we can find $h_1$, $h_0$, $h_2$ leaving the diagram commutative (three triangles and two squares to keep commutative). For this, we begin by choosing $h_0$ so as to make the middle vertical triangle commutative -- this is possible since $X_0$ is fibrant, $A_0 \longrightarrow B_0$ in $TC$.

% foot: Version 25 janvier 2004
\printedpage{36}
% head: A. Grothendieck : Dérivateurs, Ch. XIII, Édition des Universités de Montpellier II et Paris VII
\msp{page 64}

Next, we see separately for $h_1$, $h_2$, that we can find them with each satisfying the two commutativity constraints (1 vertical triangle and one tilted square). For $h_1$ e.g. existence is seen from the auxiliary square diagram (solid lines)
\[
\begin{tikzcd}
A_1 \arrow[r, "TC"] \arrow[dd] & B_1 \arrow[dd] \arrow[ddl, dotted, "h_1"'] \arrow[dr, dotted, "\alpha_1^B"] & \\
 & & B_0 \arrow[dl, dotted, "h_0"] \\
X_1 \arrow[r, "F", "\alpha_1^X"'] & X_0 \ \ . &
\end{tikzcd}
\]

\textbf{NB} This is the proof of proposition 3, page 50.

The corollary (the ``necessary'' part) having been seen, the lemma follows from it, for we know that if we have a homomorphism $f$ of cartesian squares $X$, $Y$
\[
\begin{tikzcd}[row sep=small, column sep=small]
 & X_1 \arrow[dl, "p_1^X"'] \arrow[dd, "f_1"'] & & X_3 \arrow[ll] \arrow[dl] \arrow[dd, "f_3"] \\
X_0 \arrow[dd, "f_0"'] & & X_2 \arrow[ll, crossing over, "p_2^X"'] & \\
 & Y_1 \arrow[dl, "p_1^Y"'] & & Y_3 \arrow[ll] \arrow[dl] \\
Y_0 & & Y_2 \ \ , \arrow[ll, "p_2^Y"'] \arrow[from=uu, crossing over, "f_2"] &
\end{tikzcd}
\]
with
\[
p_1^X,\ p_2^X,\ p_1^Y,\ p_2^Y \ \ \text{fibrant \texttt{[belonging to} $F$\texttt{]} (even $p_1^X$ and $p_1^Y$ fibrant is enough),}
\]
and
\[
f_0,\ f_1,\ f_2 \ \in \ W,
\]
then $f_3 \in W$. (This was seen in the theory of $W$-homotopic fiber products \texttt{[see chapter XII, p. 253]}, and the direct verification
\msp{page 65}
is immediate.) Of the relevant initial hypothesis that \begin{tikzcd}[row sep=small, column sep=small, baseline=(\tikzcdmatrixname-2-1.base)]
 & X_1 \arrow[dl] \\
X_0 & \\
 & X_2 \arrow[ul]
\end{tikzcd} and \begin{tikzcd}[row sep=small, column sep=small, baseline=(\tikzcdmatrixname-2-1.base)]
 & Y_1 \arrow[dl] \\
Y_0 & \\
 & Y_2 \arrow[ul]
\end{tikzcd}
are fibrant, we only need to use a little bit, e.g. that $p_1^X$ and $p_1^Y$ are fibrant, which ensures that the two squares $X$, $Y$ are $W$-homotopically cartesian, which is enough then to imply that $f_0$, $f_1$, $f_2 \in W \Longrightarrow f_3 \in W$.

I am going to try to generalize this reasoning, to handle the case of an ordered set $I$ which may be infinite, but such that $I^\circ$ is noetherian and sober, hence every increasing sequence of closed subsets of $I$, and every decreasing sequence of open subsets, is stationary, and moreover every upward filtering open subset has a greatest element.

% foot: Version 25 janvier 2004
\printedpage{37}
% head: A. Grothendieck : Dérivateurs, Ch. XIII, Édition des Universités de Montpellier II et Paris VII

% margin: a vertical bar is printed in the left margin beside the lines "prouver ceci : (*) Si I est tel que ... alors I aussi."
We want to prove that $I$ is admissible, and by ``noetherian induction'' we are reduced to proving this: $(*)$ If $I$ is such that for every open subset $I' \subset I$, with $I' \neq I$, $I'$ is admissible, then $I$ too.

$(*)$ granted, let us prove that $I = I_0$ is admissible.

\msp{page 66}
If it were not, by $(*)$ there would exist $I_1 \subset I_0$, $I_1 \neq I_0$, with $I_1$ not admissible. Likewise, there would be $I_2 \subset I_1$, $I_2 \neq I_1$, not admissible etc. -- we would have constructed a strictly decreasing sequence of open subsets $I_n$ of $I$, contrary to the noetherian character of $I^\circ$.

Let us prove $(*)$, and again distinguish two cases.

\textit{a}) $I^\circ$ is irreducible, hence has a generic point, \textit{i.e.} a least element, \textit{i.e.} $I$ has a greatest element. This case was already treated earlier -- every category with a final object is admissible.

\textit{b}) $I^\circ$ is reducible, hence a union of two closed subsets $I_1^\circ$, $I_2^\circ$ distinct from $I^\circ$, so $I$ is a union of two open subsets $I_1$, $I_2$ distinct from $I$. Then by the hypothesis in $(*)$, $I_1$, $I_2$ and $I_0 = I_1 \cap I_2$ are admissible. We then finish as before, being reduced to the admissibility of \begin{tikzcd}[row sep=small, column sep=small, baseline=(\tikzcdmatrixname-2-1.base)]
 & 1 & \\
\Psi = 0 \arrow[ur] \arrow[dr] & & ! \\
 & 2 &
\end{tikzcd}

\msp{page 67}
Let us prove part \textit{c}) of the theorem-scholie, part $1^\circ$). This follows from \textit{a}), once we have seen that for $s \in I$, $I'/s$ satisfies the same hypotheses as $I'$. Now $I$ being ordered, $I/s$ is an open subcategory of $I$, hence $I'/s$ an open subcategory of $I'$, with $(I'/s)^\circ$ a closed ordered subset of $I'^{\circ}$. The latter being noetherian and sober, the same holds for every closed subset, whence the conclusion.

Let us prove part $2^\circ$) of the theorem-scholie. Consider
\[
\begin{tikzcd}[row sep=large]
\cM(I')_{\rc} \arrow[r, hook, dotted] & \cM(I') & \cM(I')_{\rf} \arrow[l, hook', dotted] \arrow[d, "u_*^{\rf}"'] \\
\cM(I)_{\rc} \arrow[r, hook, dotted] \arrow[u, "u_{\rc}^*"] & \cM(I) \arrow[u, "u^*"] & \cM(I)_{\rf} \ , \arrow[l, hook', dotted]
\end{tikzcd}
\]
where $u^*$ and a fortiori $u_{\rc}^*$ are compatible with $W(I)$, $W(I')$, and by hypothesis $u_*$ is compatible with $W(I')_{\rf}$, $W(I)_{\rf}$. Passing to the categories of fractions, we find

% foot: Version 25 janvier 2004
\printedpage{38}
% head: A. Grothendieck : Dérivateurs, Ch. XIII, Édition des Universités de Montpellier II et Paris VII
\msp{page 68}
\[
\begin{tikzcd}[row sep=large]
\cM(I')_{\rc}W(I')_{\rc}^{-1} \arrow[r, dotted, "\sim"] & \HO_W(I') & \cM(I')_{\rf} \arrow[l, dotted, "\sim"'] \arrow[d, "u_*^W"'] \\
\cM(I)_{\rc}W(I)_{\rc}^{-1} \arrow[r, dotted, "\sim"] \arrow[u, "u_W^*"] & \HO_W(I) & \cM(I)_{\rf} \arrow[l, dotted, "\sim"']
\end{tikzcd}
\]
\texttt{[rather} $\cM(I')_{\rf}W(I')_{\rf}^{-1}$ \texttt{(resp.} $\cM(I)_{\rf}W(I)_{\rf}^{-1}$\texttt{) instead of} $\cM(I')_{\rf}$ \texttt{(resp.} $\cM(I)_{\rf}$\texttt{)]} and we want to show that $u_W^*$ and $u_*^W$ are adjoint, more precisely find a bifunctorial isomorphism, for $F$ \textit{in} $\cM(I)_{\rc}$ and $G'$ \textit{in} $\cM(I')_{\rf}$
\[
\Hom_W(u^*F, G') \xrightarrow{\ \sim\ } \Hom_W(F, u_*G') \tag*{$(*)$}
\]
(where $\Hom_W$ denotes the Hom in $\HO_W(I')$ resp. $\HO_W(I)$). As $F$ is in $\cM(I)_{\rc}$, $u^*F$ is in $\cM(I')_{\rc}$, and as $G'$ is in $\cM(I')_{\rf}$, $u_*G'$ is in $\cM(I)_{\rf}$. It follows that the two sides of $(*)$ are quotients of the two sides in the canonical bijection
\[
\Hom_{\cM(I')}(u^*F, G') \xrightarrow[\ \sim\ ]{\alpha} \Hom_{\cM(I)}(F, u_*G'). \tag*{$(**)$}
\]
We simply have to show that the latter transforms the relevant equivalence relation $R'$ on the first side,
\msp{page 69}
into the one $R$ on the second.

I say that in any case we have
\[
R \ \leq \ \alpha(R')
\]
\textit{without} any admissibility hypothesis on $u$, \textit{i.e.} that the inverse bijection of $(**)$ gives by passage to quotients a homomorphism % [sic] invese -> inverse, missing letter
\[
\Hom_W(u^*F, G') \longleftarrow \Hom_W(F, u_*G'), \tag*{$(***)$}
\]
necessarily surjective, and moreover functorial in $F$ \texttt{[,} $G'$ \texttt{crossed out]}, variables in $\HO_W(I)$ and $\HO_W(I')$ respectively. Indeed, we define the arrow $(*$ bis$)$ \texttt{[rather} $(***)$\texttt{]} as the composite
\[
\begin{tikzcd}[column sep=large]
\Hom_W(F, u_*G') \arrow[r, "\underbrace{\hspace{2.6cm}}_{\text{\footnotesize\begin{tabular}{@{}c@{}}since $u^*$ is\\ compatible with\\ $W(I)$, $W(I')$ on\\ the whole of $\cM(I)$,\\ and we interpret\\ the two sides\\ as Homs\\ in $\cM(I)W(I)^{-1}$\\ and $\cM(I')W(I')^{-1}$.\end{tabular}}}"'] & \Hom_W(u^*F, u^*u_*G') \arrow[dr, "\text{\footnotesize\begin{tabular}{@{}l@{}}composition\\ with the adjunction\\ homomorphism $u^*u_*G' \longrightarrow G'$\end{tabular}}"] & \\
 & & \Hom_W(u^*F, G').
\end{tikzcd}
\]
% foot: Version 25 janvier 2004

\end{document}
